\family{Fabaceae}{Pea family}
\label{fam:fabaceae}

Trees, shrubs, climbers or herbs. Leaves alternate, usually pinnate or
bipinnate, with stipules and with a swollen pulvinus at the base of
the petiole. Flowers regular or irregular, often showy. Fruit a pod
(legume), sometimes indehiscent and winged. Three species in the Arami
Uplands, all trees of the forest--savanna mosaic and savanna.

\section*{Key to the species}
\begin{florakey}
	\lead{1a}{Leaves bipinnate; flowers in pendent red heads on long
		stalks}{\taxon{Parkia biglobosa}}
	\lead{1b}{Leaves once pinnate}{2}
	\lead{2a}{Leaflets in 3--5 opposite pairs, without a terminal leaflet;
		fruit a thick woody pod; seeds black with an orange
		aril}{\taxon{Afzelia africana}}
	\lead{2b}{Leaflets alternate, with a terminal leaflet; fruit a flat,
		round, winged disc, bristly over the seed}{\taxon{Pterocarpus
			erinaceus}}
\end{florakey}

\species{Afzelia africana}{Sm. ex Pers.}
\label{sp:FAB01}

\field{Local and trade names} Yoruba: Apa; Hausa: Kawo. Trade:
Doussie.

\field{Status} Native.

\field{Description} Deciduous tree to 30~m, smaller in savanna, with a
spreading crown. Bark grey, flaking in patches. Leaves paripinnate
with 3--5 pairs of opposite, ovate-elliptic leaflets 5--10~cm long.
Flowers with a single large petal, white streaked with red. Pod woody,
flattened, 12--17~cm long, black when ripe. Seeds black, each with a
bright orange aril.

\field{Phenology} Flowering March to May; fruiting November to
February.

\field{Habitat} Forest--savanna mosaic and Guinea savanna woodland.

\field{Uses} A heavy, durable timber used for flooring, joinery and
boat-building. The leaves are lopped for livestock fodder
\citep{arbonnier2004,louppe2008}.

\field{Specimens examined} Agbede (C2), 540~m, st., 13~Sep 1936, \specimen{Harding}{1273}{OKF};
Agbede (C2), 540~m, fr., 13~Nov 1936, \specimen{Harding}{1296}{OKF};
Ewe Stream gallery forest (B3), 580~m, fr., 6~Feb 1965, \specimen{Olaniyan}{577}{SVN};
Aramiwo Escarpment, summit (B1), 900~m, st., 25~Oct 1965, \specimen{Olaniyan}{663}{SVN};
Aramiwo Escarpment, summit (B1), 900~m, st., 5~Aug 2022, \specimen{Alade}{493}{AUH}.

\field{Distribution} Grid squares B1, B3, C2.

\begin{center}
	\distmap{B/1, B/3, C/2}
\end{center}

\species{Parkia biglobosa}{(Jacq.) R.Br. ex G.Don}
\label{sp:FAB02}

\field{Local and trade names} Yoruba: Igba; Hausa: Dorowa. Trade:
African locust bean.

\field{Status} Native.

\field{Description} Deciduous tree to 20~m with a wide, spreading
crown. Bark thick, dark grey-brown, fissured. Leaves bipinnate with
many pairs of pinnae, each bearing numerous small leaflets. Flowers in
dense, pendent, red spherical heads on long stalks. Pods in clusters,
20--30~cm long, brown, the seeds embedded in a yellow, mealy pulp.

\field{Phenology} Flowering December to February; fruiting March to
May.

\field{Habitat} Forest--savanna mosaic and Guinea savanna woodland;
usually kept when land is cleared for farming.

\field{Uses} The fermented seeds are made into a widely used condiment
(Yoruba \textit{iru}, Hausa \textit{daddawa}). The sweet yellow pulp
is eaten \citep{arbonnier2004,orwa2009}.

\field{Specimens examined} Agbede (C2), 540~m, st., 18~Nov 1937, \specimen{Harding}{1473}{OKF};
Agbede (C2), 540~m, fr., 9~Mar 1938, \specimen{Harding}{1628}{OKF};
Asiri Crossing (A2), 740~m, fl., 20~Jan 1965, \specimen{Olaniyan}{515}{SVN};
Asiri Crossing (A2), 740~m, fr., 7~Apr 1967, \specimen{Olaniyan}{872}{SVN};
Okuta Hill (A4), 820~m, fr., 1~May 2022, \specimen{Alade}{476}{AUH}.

\field{Distribution} Grid squares A2, A4, C2.

\begin{center}
	\distmap{A/2, A/4, C/2}
\end{center}

\species{Pterocarpus erinaceus}{Poir.}
\label{sp:FAB03}

\field{Local and trade names} Hausa: Madobiya. Trade: African
rosewood.

\field{Status} Native.

\field{Description} Deciduous tree to 15~m. Bark dark grey and scaly,
exuding a red gum when cut. Leaves imparipinnate with alternate
leaflets. Flowers yellow, in panicles, appearing while the tree is
leafless. Fruit a flat, round, winged disc 4--7~cm across, bristly
over the seed.

\field{Phenology} Flowering December to February; fruiting February
to April.

\field{Habitat} Guinea savanna woodland.

\field{Uses} The red-brown timber is used for furniture and carving,
and the leaves are lopped for fodder. Heavy logging for export has
made it scarce in much of its range \citep{arbonnier2004,orwa2009}.

\field{Specimens examined} Asiri Crossing (A2), 740~m, fr., 15~Apr 1964, \specimen{Olaniyan}{467}{SVN};
Asiri Crossing (A2), 740~m, fr., 14~Apr 1966, \specimen{Olaniyan}{719}{SVN};
Idofin (A3), 760~m, fl., 4~Dec 1967, \specimen{Olaniyan}{912}{SVN};
Asiri Crossing (A2), 740~m, fl., 6~Dec 2020, \specimen{Alade}{288}{AUH};
Idofin (A3), 760~m, fr., 27~Feb 2024, \specimen{Alade}{623}{AUH}.

\field{Distribution} Grid squares A2, A3.

\begin{center}
	\distmap{A/2, A/3}
\end{center}